% AUTO-GENERATED by figs/make_figures.py -- do not edit.
\begin{figure}[t]\centering
\begin{tikzpicture}
\begin{axis}[paperfig,width=\columnwidth,height=0.66\columnwidth,
  xlabel={replication degree $k$},ylabel={mean $\PAoIper$ (normalized)},
  legend pos=north west,grid=both,legend cell align=left,xtick=data]
\addplot[thick,cA,mark=*,mark size=1.3pt] coordinates {(1.0,1.410314663389778) (2.0,1.248192570863487) (3.0,1.1777586195908003) (4.0,1.4564704134531814) (5.0,1.7592147238268279) (6.0,2.073176697683509) (7.0,2.3922001335971532) (8.0,2.7185592470423137) (9.0,3.0510884750494167)};
\addlegendentry{$\eta=3$ ($k^\ast=3$)}
\addplot[thick,cB,mark=square*,mark size=1.3pt] coordinates {(1.0,1.4101859347850572) (2.0,1.2478719880669222) (3.0,1.1692371885105677) (4.0,1.123112635056019) (5.0,1.1020916179657483) (6.0,1.2732102430111434) (7.0,1.458191515479616) (8.0,1.6476294592829408) (9.0,1.8487261126269243)};
\addlegendentry{$\eta=5$ ($k^\ast=5$)}
\end{axis}
\end{tikzpicture}
\caption{Result~3 (E7): the integrated energy+replication+AoI objective
is unimodal in $k$ with minimum at one of the two integers bracketing
$\eta$; pushing
$k$ past the energy budget starves future contacts and raises the peak.}
\label{fig:r3unimodal}
\end{figure}
